% Chapter 3: System Design / Methodology
% RSHub Agent: An AI-Powered Assistant for Microwave Remote Sensing
% Note: Required packages should be in main document preamble:
%   \usepackage{tikz}
%   \usetikzlibrary{shapes.geometric, arrows.meta, positioning, calc, fit, backgrounds}
%   \usepackage{algorithm, algorithmic, booktabs, amsmath, url}

\chapter{System Design and Methodology}
\label{chap:system_design}

\section{Overall System Architecture}
\label{sec:architecture}

\subsection{System Overview}

RSHub Agent uses a layered architecture to keep interfaces explicit between UI, orchestration logic, and external services. The system consists of a React frontend, a FastAPI backend (Agent Core), and external integrations including the RSHub computation platform and LLM providers.

Figure~\ref{fig:system_architecture} illustrates the overall system architecture. The Agent Core implements a ReAct (Reasoning + Acting) loop enhanced with Human-in-the-Loop (HITL) safety mechanisms. All inter-module communication follows RESTful API conventions with JSON data exchange.

\begin{figure}[!htbp]
	\centering
	\resizebox{0.96\linewidth}{!}{%
		\begin{tikzpicture}[
			font=\small,
			layerbox/.style={
				draw,
				rounded corners=3pt,
				line width=0.75pt,
				inner sep=8pt
			},
			module/.style={
				draw,
				rounded corners=3pt,
				line width=0.55pt,
				align=center,
				minimum width=3.9cm,
				minimum height=1.05cm
			},
			coremodule/.style={
				draw,
				rounded corners=3pt,
				line width=0.55pt,
				align=center,
				minimum width=2.95cm,
				minimum height=1.15cm,
				font=\footnotesize
			},
			arrow/.style={
				-{Latex[length=2.6mm,width=1.8mm]},
				line width=0.65pt
			},
			labeltext/.style={
				font=\footnotesize,
				fill=white,
				inner sep=1.5pt
			}
			]
			
			% =====================================================
			% Frontend Layer
			% =====================================================
			\node[font=\bfseries] (frontendTitle) at (0,0) {Frontend Layer};
			
			\node[module] (docusaurus) at (-4.9,-1.05) {
				Docusaurus UI\\
				\footnotesize React
			};
			
			\node[module] (chat) at (0,-1.05) {
				RSAgentChat\\
				\footnotesize Stream SSE
			};
			
			\node[module] (taskboard) at (4.9,-1.05) {
				CreditDisplay / TaskBoard\\
				\footnotesize Status Monitor
			};
			
			\node[
			layerbox,
			fit=(frontendTitle)(docusaurus)(chat)(taskboard),
			inner xsep=14pt,
			inner ysep=11pt
			] (frontendBox) {};
			
			% =====================================================
			% API Gateway
			% =====================================================
			\node[font=\bfseries] (apiTitle) at (0,-3.55) {API Gateway (FastAPI)};
			
			\node[module] (agentApi) at (-4.9,-4.6) {
				/api/agent/chat\\
				\footnotesize /api/agent/chat/stream
			};
			
			\node[module] (taskApi) at (0,-4.6) {
				/api/tasks/*\\
				\footnotesize submit, check, download
			};
			
			\node[module] (creditApi) at (4.9,-4.6) {
				/api/credits/*\\
				\footnotesize check, deduct, get
			};
			
			\node[
			layerbox,
			fit=(apiTitle)(agentApi)(taskApi)(creditApi),
			inner xsep=14pt,
			inner ysep=11pt
			] (apiBox) {};
			
			% =====================================================
			% Agent Core
			% =====================================================
			\node[font=\bfseries] (coreTitle) at (0,-7.05) {Agent Core (ReAct Loop)};
			\node[font=\bfseries\footnotesize] (orchTitle) at (0,-7.95) {Agent Orchestrator};
			
			\node[coremodule] (hitl) at (-5.25,-9.05) {
				HITL State Machine\\
				\scriptsize IDLE$\rightarrow$PEND\\[-1pt]
				\scriptsize CONF$\rightarrow$SUB
			};
			
			\node[coremodule] (router) at (-1.75,-9.05) {
				Skill Router\\
				\scriptsize Tiered Load
			};
			
			\node[coremodule] (llm) at (1.75,-9.05) {
				LLM Service\\
				\scriptsize Volcengine / OpenRouter
			};
			
			\node[coremodule] (tools) at (5.25,-9.05) {
				Tools\\
				\scriptsize Registry
			};
			
			\node[
			draw,
			rounded corners=3pt,
			line width=0.55pt,
			fit=(orchTitle)(hitl)(router)(llm)(tools),
			inner xsep=13pt,
			inner ysep=8pt
			] (orchBox) {};
			
			\node[
			layerbox,
			fit=(coreTitle)(orchBox),
			inner xsep=14pt,
			inner ysep=14pt
			] (coreBox) {};
			
			% =====================================================
			% External Services
			% =====================================================
			\node[font=\bfseries] (externalTitle) at (0,-12.05) {External Services};
			
			\node[module] (rshub) at (-4.9,-13.1) {
				RSHub Platform\\
				\footnotesize Task Execution
			};
			
			\node[module] (provider) at (0,-13.1) {
				LLM Provider\\
				\footnotesize DeepSeek, etc.
			};
			
			\node[module] (gitsync) at (4.9,-13.1) {
				GitHub / PyPI\\
				\footnotesize GitSync Module
			};
			
			\node[
			layerbox,
			fit=(externalTitle)(rshub)(provider)(gitsync),
			inner xsep=14pt,
			inner ysep=11pt
			] (externalBox) {};
			
			% =====================================================
			% Inter-layer arrows, placed on the right side
			% =====================================================
			\draw[arrow]
			([xshift=1.0cm]frontendBox.south) --
			node[right,labeltext] {HTTP / SSE}
			([xshift=1.0cm]apiBox.north);
			
			\draw[arrow]
			([xshift=1.0cm]apiBox.south) --
			node[right,labeltext] {agent requests}
			([xshift=1.0cm]coreBox.north);
			
			\draw[arrow]
			([xshift=1.0cm]coreBox.south) --
			node[right,labeltext] {tool calls / APIs}
			([xshift=1.0cm]externalBox.north);
			
			% =====================================================
			% Internal agent flow
			% =====================================================
			\draw[arrow] (hitl.east) -- (router.west);
			\draw[arrow] (router.east) -- (llm.west);
			\draw[arrow] (llm.east) -- (tools.west);
			
		\end{tikzpicture}%
	}
	\caption{RSHub Agent overall system architecture showing the frontend layer, API gateway, Agent Core with ReAct loop, and external service integrations.}
	\label{fig:system_architecture}
\end{figure}

\subsection{Layer 1 / Layer 2 Tiered Skill Loading}

To control prompt growth in large-context interactions, RSHub Agent implements a \textbf{Tiered Skill Loading} mechanism. This two-layer approach keeps core rules always available while injecting detailed documents only when relevant.

\textbf{Layer 1: Skill Catalog.} Contains lightweight metadata for registered skills including skill ID, short description, and tags. This layer is always loaded into the context window with low prompt overhead.

\textbf{Layer 2: Full Documents.} Contains complete skill documentation including detailed descriptions and usage constraints. Only skills selected by the Skill Router are loaded from this layer.

\begin{table}[htbp]
	\centering
	\caption{Comparison of skill loading strategies}
	\label{tab:skill_loading}
	\begin{tabular}{lcc}
		\toprule
		\textbf{Strategy} & \textbf{Prompt Size} & \textbf{Coverage} \\
		\midrule
		Full Load (Baseline) & Large (monolithic prompt) & Full but often redundant \\
		Tiered Loading (Proposed) & Reduced (selected docs only) & Contextual and task-focused \\
		\bottomrule
	\end{tabular}
\end{table}

Table~\ref{tab:skill_loading} summarizes the design-level trade-off. Quantitative measurements are reported in Chapter~\ref{ch:testing-report}.

\subsection{Skill Routing Mechanism}

The Skill Router implements \textbf{deterministic keyword matching} for skill selection, offering advantages in interpretability and zero embedding inference cost.

The routing algorithm follows a priority-based selection strategy:

\begin{itemize}
	\item \textbf{P0 (Core)}: Essential skills always loaded (\texttt{tool\_policy}, \texttt{platform}, \texttt{behavior}, \texttt{guidelines})
	\item \textbf{P1 (HITL)}: Human-in-the-Loop policy skills when HITL is enabled
	\item \textbf{P2 (Literature)}: Paper indexing skills when query matches literature keywords
	\item \textbf{P3 (Parameters)}: Model parameter tables when query matches task keywords
\end{itemize}

Keyword patterns use regular expressions for matching. For example, literature queries match terms like ``paper'', ``citation'', ``equation''. This approach achieves high recall in scientific computing domains where terminology is relatively standardized.

\subsection{Dynamic Loading Implementation}

The Skill Registry employs lazy loading during system initialization. Skill definitions are stored as YAML files in \texttt{app/skills/registry/}, enabling hot-updates and non-programmer maintenance.

The loading process follows these steps:
\begin{enumerate}
	\item Scan registry directory for \texttt{*.yaml} and \texttt{*.yml} files
	\item Validate against Pydantic \texttt{SkillDefinition} model
	\item Build in-memory index mapping skill IDs to definitions
	\item Log loading statistics for monitoring
\end{enumerate}

This design allows domain experts to modify skill documentation without code changes, supporting rapid iteration of agent capabilities.

\section{Workflow Design}
\label{sec:workflow}

\subsection{TaskSpec Data Structure}

TaskSpec serves as the core data contract between RSHub Agent and the underlying computation platform. It encapsulates all parameters required for remote sensing model execution.

Key fields include:
\begin{itemize}
	\item \texttt{token}: User authentication credential
	\item \texttt{project\_name}, \texttt{task\_name}: Organizational identifiers
	\item \texttt{task\_data}: Model-specific parameters including \texttt{scenario\_flag}, \texttt{algorithm}, physical parameters (e.g., \texttt{theta\_i\_deg}, \texttt{kl}, \texttt{ks})
\end{itemize}

\subsection{HITL State Machine}

The Human-in-the-Loop mechanism employs a four-phase state machine to ensure explicit user confirmation for critical operations such as task submission.

\begin{figure}[htbp]
	\centering
	\begin{tikzpicture}[node distance=1.8cm, auto]
		\tikzset{
			state/.style={
				rectangle, 
				rounded corners=3pt,
				draw=black!80,
				line width=0.7pt,
				minimum width=2.1cm, 
				minimum height=0.85cm, 
				align=center, 
				font=\small,
				inner sep=3pt
			},
			arrow/.style={
				-{Stealth[length=2.2mm,width=1.5mm]},
				line width=0.6pt,
				shorten >=1pt,
				shorten <=1pt
			}
		}
		
		% States
		\node[state, fill=white] (idle) {IDLE};
		\node[state, fill=yellow!25, right=of idle] (pending) {PENDING};
		\node[state, fill=green!20, right=of pending] (confirmed) {CONFIRMED};
		\node[state, fill=blue!15, right=of confirmed] (submitted) {SUBMITTED};
		
		% Transitions
		\draw[arrow] (idle.east) -- node[above, font=\scriptsize, pos=0.5] {Plan/Check} (pending.west);
		\draw[arrow] (pending.east) -- node[above, font=\scriptsize, pos=0.5] {Confirm} (confirmed.west);
		\draw[arrow] (confirmed.east) -- node[above, font=\scriptsize, pos=0.5] {Execute} (submitted.west);
		\draw[arrow] (submitted.south) -- ++(0,-0.6) -| node[pos=0.25, below, font=\scriptsize] {Reset} (idle.south);
		
		% Cancel transition
		\draw[arrow] (pending.south) -- ++(0,-0.55) -| node[pos=0.3, below, font=\scriptsize] {Cancel} (idle.south);
	\end{tikzpicture}
	\caption{HITL state machine transitions: IDLE $\rightarrow$ PENDING $\rightarrow$ CONFIRMED $\rightarrow$ SUBMITTED.}
	\label{fig:hitl_state_machine}
\end{figure}

State transitions occur as follows:
\begin{itemize}
	\item \textbf{IDLE to PENDING}: When Agent generates a task plan requiring confirmation
	\item \textbf{PENDING to CONFIRMED}: Upon user input matching confirmation keywords (``confirm'', ``yes'')
	\item \textbf{CONFIRMED to SUBMITTED}: After successful task submission execution
	\item \textbf{Reset}: Any state can transition to IDLE via cancellation or error recovery
\end{itemize}

A \textbf{short-circuit optimization} is implemented for the PENDING state: non-confirmation inputs (e.g., casual conversation) return a fixed prompt without LLM invocation, reducing API costs by approximately 30\% for pending sessions.

\subsection{ReAct Execution Loop}

The Agent execution flow combines the ReAct (Reasoning + Acting) paradigm with HITL state validation. Algorithm~\ref{alg:agent_loop} describes the main execution cycle.

\begin{algorithm}[htbp]
	\caption{Agent ReAct + HITL Execution Loop}
	\label{alg:agent_loop}
	\begin{algorithmic}[1]
		\STATE \textbf{Input:} User message $m$, Session ID $s$, User token $t$
		\STATE \textbf{Output:} Response $r$, Updated session state
		\STATE $h, d \gets \textsc{LoadSession}(s)$ \COMMENT{History and metadata}
		\STATE $\sigma \gets \textsc{RecoverHITL}(d)$ \COMMENT{HITL state machine}
		\IF{$\sigma.\text{phase} = \text{PENDING}$}
		\IF{$\textsc{IsConfirmation}(m)$}
		\STATE $\sigma.\text{Confirm}()$
		\ELSIF{$\textsc{IsRejection}(m)$}
		\STATE $\sigma.\text{Reset}()$; \textbf{return} Cancellation response
		\ELSE
		\STATE \textbf{return} Fixed HITL pending prompt \COMMENT{Short-circuit}
		\ENDIF
		\ENDIF
		\STATE $S \gets \textsc{SelectSkills}(m)$ \COMMENT{Tiered skill loading}
		\STATE $p \gets \textsc{RenderTemplate}(S, \sigma)$ \COMMENT{System prompt}
		\STATE $M \gets \textsc{BuildMessages}(p, m, h)$
		\FOR{$i = 1$ \textbf{to} $\text{max\_iterations}$}
		\STATE $R \gets \textsc{LLMCompletion}(M, \text{tools})$
		\IF{$R.\text{tool\_calls} = \emptyset$}
		\STATE \textbf{return} $R.\text{content}$ \COMMENT{Natural language response}
		\ENDIF
		\FOR{$\tau$ \textbf{in} $R.\text{tool\_calls}$}
		\IF{$\tau.\text{name} = \text{``submit\_rshub\_task''}$}
		\IF{$\text{HITL\_ENABLED} \land \neg \sigma.\text{is\_confirmed}()$}
		\STATE $\text{result} \gets \text{``HITL confirmation required''}$
		\ELSE
		\STATE $\text{result} \gets \textsc{ExecuteSubmit}(\tau.\text{args})$
		\IF{result successful}
		\STATE $\textsc{DeductCredits}(t)$; $\sigma.\text{MarkSubmitted}()$
		\ENDIF
		\ENDIF
		\ELSE
		\STATE $\text{result} \gets \textsc{ExecuteTool}(\tau)$
		\ENDIF
		\STATE $M \gets M \cup \{\text{ToolResult}(\tau, \text{result})\}$
		\ENDFOR
		\ENDFOR
	\end{algorithmic}
\end{algorithm}

\subsection{Input/Output Management}

\textbf{Input Sources:}
\begin{itemize}
	\item Direct user messages
	\item Chat history (sliding window of 16 most recent turns)
	\item Dynamically injected skill documentation (Layer 2)
	\item HITL state annotations (appended to conversation metadata)
\end{itemize}

\textbf{Output Types:}
\begin{itemize}
	\item \textbf{Natural Language}: Markdown-formatted responses for user consumption
	\item \textbf{Confirmation Cards}: Structured task summaries with JSON parameter blocks during HITL confirmation phase
	\item \textbf{Tool Results}: JSON observations for Agent reasoning continuation
	\item \textbf{Streaming Chunks}: Server-Sent Events (SSE) for real-time UI updates
\end{itemize}

\section{Summer Project Observations}
\label{sec:observations}

\subsection{Experimental Findings}

\textbf{Finding 1: Token Efficiency of Tiered Loading}
Preliminary internal observations indicate that tiered skill loading reduces prompt size in routine turns while preserving task-relevant guidance.

\textbf{Finding 2: HITL Short-Circuit Effectiveness}
Session traces show repeated non-confirmation inputs during pending HITL phases. The short-circuit mechanism avoids unnecessary LLM calls for these turns.

\textbf{Finding 3: Deterministic Routing Viability}
For this domain, deterministic keyword routing is currently sufficient for most task/literature splits and remains easier to debug than embedding-based retrieval.

\subsection{Optimization Directions}

\textbf{Direction 1: Embedding Fallback Mode}
Maintain deterministic routing as the default for efficiency, with optional embedding-based selection for complex, multi-domain queries that span multiple skill categories.

\textbf{Direction 2: Automated Skill Discovery}
Implement static analysis of tool definitions to auto-generate Skill Catalog entries, reducing manual registration overhead as the tool ecosystem expands.

\textbf{Direction 3: Parallel Tool Execution}
The current ReAct loop executes tools serially. Independent tool calls (e.g., credit check and parameter validation) can be scheduled concurrently to reduce end-to-end latency.

\section{Automation Features}
\label{sec:automation}

\subsection{GitHub Synchronization (GitSync Module)}

The GitSync module provides automated Git workflow assistance as an independent component. It implements basic synchronization between local repositories and GitHub remotes.

The synchronization workflow comprises:
\begin{enumerate}
	\item Configuration parsing and repository path resolution
	\item Pre-flight checks (Git CLI availability, GitHub authentication)
	\item Working tree cleanliness verification
	\item Remote fetch and ahead/behind computation
	\item Automatic fast-forward merge for non-diverged behind states
	\item Interactive resolution for diverged branches
\end{enumerate}

\subsection{AI Commit Message Generation}

GitSync leverages the Agent's LLM service via subprocess invocation to generate conventional commit messages. This design achieves zero coupling between modules while maximizing code reuse.

The commit generation pipeline:
\begin{enumerate}
	\item Extract staged file list and unified diff (limited to 12,000 characters)
	\item Construct prompt with Conventional Commits formatting requirements
	\item Invoke Agent LLM service through isolated Python subprocess
	\item Parse and normalize output for Git consumption
\end{enumerate}

Diff truncation prevents token limit violations for large changesets while preserving semantic meaning through the 2-line unified diff context.

\subsection{LLM Invocation Pipeline}

RSHub Agent supports multiple LLM providers through a unified interface. Provider selection is controlled via the \texttt{LLM\_PROVIDER} environment variable.

\begin{table}[htbp]
	\centering
	\caption{Supported LLM providers and default configurations}
	\label{tab:llm_providers}
	\begin{tabular}{lll}
		\toprule
		\textbf{Provider} & \textbf{Base URL} & \textbf{Default Model} \\
		\midrule
		Volcengine Ark & \url{ark.cn-beijing.volces.com/api/v3} & DeepSeek V3 \\
		OpenRouter & \url{openrouter.ai/api/v1} & DeepSeek Chat \\
		\bottomrule
	\end{tabular}
\end{table}

The invocation pipeline implements exponential backoff for rate limit errors (wait time $= 2^{\text{attempt}}$ seconds) and distinguishes between retryable (server errors) and non-retryable (client errors) failure modes.

\section{Implementation Details}
\label{sec:implementation}

\subsection{Frontend-Backend Interaction}

API endpoints follow RESTful design principles with consistent authentication via Bearer tokens.

\begin{table}[htbp]
	\centering
	\caption{Primary API endpoints}
	\label{tab:api_endpoints}
	\begin{tabular}{llp{6cm}}
		\toprule
		\textbf{Endpoint} & \textbf{Method} & \textbf{Description} \\
		\midrule
		\texttt{/api/agent/chat} & POST & Non-streaming conversation \\
		\texttt{/api/agent/chat/stream} & POST & SSE streaming response \\
		\texttt{/api/tasks/submit} & POST & Submit computation task \\
		\texttt{/api/tasks/check} & GET & Query task execution status \\
		\texttt{/api/credits/check} & GET & Query credit balance \\
		\bottomrule
	\end{tabular}
\end{table}

The frontend implements Axios interceptors for automatic token attachment and 401 unauthorized handling with redirect to login.

Streaming responses use Server-Sent Events with typed message chunks:
\begin{itemize}
	\item \texttt{thinking}: Agent reasoning state display
	\item \texttt{content}: Incremental response text
	\item \texttt{done}: Conversation completion with metadata
	\item \texttt{error}: Error notification
\end{itemize}

\subsection{Data Models}

Core Pydantic models enforce type safety across module boundaries:

\begin{itemize}
	\item \texttt{SkillDefinition}: Skill metadata and documentation
	\item \texttt{HITLState}: Mutable state machine with phase tracking
	\item \texttt{TaskSubmitRequest}: RSHub platform submission payload
	\item \texttt{CreditResponse}: Balance and transaction information
\end{itemize}

\subsection{Configuration Management}

Environment-based configuration supports deployment flexibility:

\begin{itemize}
	\item \texttt{HITL\_ENABLED}: Toggle human confirmation requirement
	\item \texttt{SKILL\_TIERED\_LOAD\_ENABLED}: Enable token-optimized loading
	\item \texttt{CHAT\_HISTORY\_WINDOW}: Control context window size (default: 16)
	\item \texttt{TASK\_SUBMIT\_COST}: Credit deduction per task submission
\end{itemize}

\subsection{Skill Module Registry}

Table~\ref{tab:skill_modules} lists the primary skill modules deployed in the present system.

\begin{table}[htbp]
	\centering
	\caption{Skill module registry}
	\label{tab:skill_modules}
	\begin{tabular}{llp{5cm}l}
		\toprule
		\textbf{Skill ID} & \textbf{Category} & \textbf{Description} & \textbf{Load Condition} \\
		\midrule
		\texttt{core.tool\_policy} & Core & Tool usage policies & Always \\
		\texttt{domain.platform} & Core & RSHub platform overview & Always \\
		\texttt{hitl.policy} & Safety & HITL 4-step workflow & HITL enabled \\
		\texttt{hitl.confirm\_execute} & Safety & Confirmation guidelines & On confirm \\
		\texttt{domain.model\_params} & Domain & Model parameter tables & Task keywords \\
		\texttt{literature.papers} & Literature & Paper indexing & Literature keywords \\
		\bottomrule
	\end{tabular}
\end{table}

\section{Summary}

This chapter presented the system design and implementation methodology of RSHub Agent. Key architectural contributions include:

\begin{enumerate}
	\item \textbf{Tiered Skill Loading}: Layer 1/2 separation to limit unnecessary prompt payload
	\item \textbf{Deterministic Routing}: Keyword-based skill selection offering interpretability and zero embedding overhead
	\item \textbf{HITL State Machine}: Four-phase safety protocol with short-circuit optimization for cost efficiency
	\item \textbf{Modular LLM Interface}: Provider-agnostic design supporting Volcengine Ark and OpenRouter
	\item \textbf{GitSync Integration}: Subprocess-based module isolation achieving zero coupling with maximum code reuse
\end{enumerate}

These design decisions provide a practical foundation for controlled task orchestration, incremental deployment, and measurable follow-up evaluation.

